\section{The exact finite dihedral character identity}
\label{sec:cubic-finite-type-character}

The following finite-group calculation supplies the character input in
AC-G72.  It makes no assertion that the necessary local automorphic
constituent occurs.  Let $q$ be an odd prime, let $p>3$ divide $q+1$,
and let $\mathcal O$ be a sufficiently large coefficient discrete
valuation ring of residue characteristic $p$, with fraction field $E$
and residue field $k$.  Enlarge $E$ finitely whenever necessary to
split the finite groups below.  Write
\[
 G_2=\operatorname{GL}_2(\mathbb F_q),\qquad
 G_3=\operatorname{GL}_3(\mathbb F_q),\qquad
 T'=\mathbb F_{q^2}^{\times}/\mathbb F_q^{\times}.
                                                        \tag{G72C.1}
\]
Choose a nontrivial $p$-power-order character $\theta$ of $T'$.  Then
$\theta^q=\theta^{-1}$ and $\theta^2\ne1$.  Let $P\subset G_3$
stabilize the plane spanned by the first two standard basis vectors;
its Levi quotient is $G_2\times\mathbb F_q^{\times}$.

\begin{proposition}[AC-G72-character: a finite dihedral type]
\label{AC-G72-character}
There is an irreducible cuspidal representation $C_\theta$ of $G_2$,
with trivial central character and dimension $q-1$, for which
\[
 \tau_\theta=\operatorname{Ind}_{P}^{G_3}(C_\theta\boxtimes1)
                                                        \tag{G72C.2}
\]
is irreducible of dimension $q^3-1$.  Here the representation of the
Levi is inflated to $P$.  For arbitrary stable $\mathcal O$-lattices,
reduction gives the identities in the Grothendieck groups of finite
$k$-representations
\[
 [\overline C_\theta]=[\overline{\operatorname{St}}_2]-[1],
 \qquad
 [\overline\tau_\theta]=[\overline{\operatorname{St}}_3]-[1].
                                                        \tag{G72C.3}
\]
Equivalently, the corresponding ordinary characters agree on
$p$-regular elements.  The characteristic-zero identities needed
for the second formula are
\[
 \operatorname{Ind}_{P}^{G_3}(1\boxtimes1)=1+\rho_{21},
 \qquad
 \operatorname{Ind}_{P}^{G_3}(\operatorname{St}_2\boxtimes1)
                     =\rho_{21}+\operatorname{St}_3,     \tag{G72C.4}
\]
where $\rho_{21}$ is irreducible of dimension $q^2+q$.
\end{proposition}
\begin{proof}
\proofstep{An explicit rank-two character.}
Work first in $\overline G_2=\operatorname{PGL}_2(\mathbb F_q)$.
Let $U$ be its standard upper unipotent group and choose a
nontrivial character $\psi$ of $U$.  Consider the virtual character
\[
       c_\theta=\operatorname{Ind}_{U}^{\overline G_2}\psi
                         -\operatorname{Ind}_{T'}^{\overline G_2}\theta.
                                                        \tag{G72C.5}
\]
This is the explicit formula in Lang--Pollack--Wake,
\emph{Counting level-raising congruences using modular representation
theory}, arXiv:2609.05144v1, Section~2.2, printed p.~10.  We verify
all its properties used here directly.  The group has order
$q(q^2-1)$.  Its nonidentity unipotent elements form one conjugacy
class with centralizer of order $q$.  Noncentral split semisimple
elements meet neither $U$ nor $T'$.  A nonsplit semisimple class
meets $T'$ in $\{t,t^{-1}\}$; its centralizer has order $q+1$,
except for the involution, where the order is $2(q+1)$.
The induced-character formula therefore gives
\[
\begin{array}{c|cccc}
 &1&\text{nonidentity unipotent}&\text{split regular}
                                      &\text{nonsplit }t\ne1\\ \hline
 c_\theta&q-1&-1&0&-\theta(t)-\theta(t)^{-1}.
\end{array}                                                   \tag{G72C.6}
\]
For the unipotent value we used
$\sum_{u\in U\setminus\{1\}}\psi(u)=-1$.
The last formula also holds for the nonsplit involution: the
single torus intersection and the doubled centralizer contribute
$2\theta(t)$ before the minus sign.

As $\theta^2\ne1$, summing its square and inverse square over $T'$
gives zero.  Consequently the ordinary character norm is
\[
 \langle c_\theta,c_\theta\rangle
 =\frac{(q-1)^2}{q(q^2-1)}+\frac1q
   +\frac{1}{2(q+1)}\sum_{t\in T'\setminus\{1\}}
                          |\theta(t)+\theta(t)^{-1}|^2
 =\frac{q-1}{q(q+1)}+\frac1q+1-\frac2{q+1}=1.
                                                        \tag{G72C.7}
\]
The coefficient $1/(2(q+1))$ handles both ordinary inverse pairs
and the involution correctly.  A virtual character of norm one
is plus or minus one irreducible character, because its integral
irreducible-character coefficients have squares summing to one.
The positive degree $q-1$ selects the plus sign.  Moreover
\[
 \dim c_\theta^U
       =q^{-1}\bigl((q-1)+(q-1)(-1)\bigr)=0.             \tag{G72C.8}
\]
Inflation to $G_2$ gives the asserted cuspidal $C_\theta$, with
trivial scalar action.  For rank two, vanishing of invariants
under the unipotent radical of the proper parabolic is precisely
the cuspidality condition.

The rank-two Steinberg character is the permutation character
on $\mathbb P^1(\mathbb F_q)$ minus the trivial character.  Its
values in the four columns of (G72C.6) are $q,0,1,-1$.
On $p$-regular nonsplit elements the character $\theta$ equals
one, since it has $p$-power order.  Thus $c_\theta$ and
$\operatorname{St}_2-1$ agree on every $p$-regular element.
The same is true after inflation to $G_2$.  The injectivity of
Brauer characters on the Grothendieck group gives the first
identity in (G72C.3); it is independent of the chosen lattices.

\proofstep{The finite flag calculation in rank three.}
Let $B\subset G_3$ be the standard Borel, and put
$A=\operatorname{Ind}_{B}^{G_3}1$ and
$M=\operatorname{Ind}_{P}^{G_3}1$.
The action on planes has two orbits on ordered pairs, according
as the planes coincide or intersect in a line.  Thus
$\langle M,M\rangle=2$ and $M=1+\rho_{21}$, with
$\rho_{21}$ irreducible.  Full flags have six relative positions,
and a plane has three relative positions with respect to a full
flag.  These elementary relative-position counts give
\[
 \langle A,A\rangle=6,\quad
 \langle A,M\rangle=3,\quad
 \langle A,1\rangle=\langle M,1\rangle=1.                 \tag{G72C.9}
\]
The virtual character $A-2M+1$ has norm
$6+8+1-12+2-4=1$ and degree
\[
 (q^2+q+1)(q+1)-2(q^2+q+1)+1=q^3.
\]
It is therefore irreducible.  The line and plane permutation
characters are equal: their incidence matrix has product with
its transpose equal to $qI+J$, where $J$ is the all-ones matrix,
and is consequently an isomorphism in characteristic zero.
Thus $A-2M+1$ is the usual alternating building character,
namely the rank-three Steinberg character.  Transitivity of induction and
$\operatorname{Ind}_{B\cap G_2}^{G_2}1=1+\operatorname{St}_2$
give
$A=M+\operatorname{Ind}_{P}^{G_3}(\operatorname{St}_2\boxtimes1)$.
Combining this with $A-2M+1=\operatorname{St}_3$ proves (G72C.4).
Induction and inflation are exact over $k$, so applying them
to the first identity of (G72C.3) proves its second identity.

\proofstep{Irreducibility of the induced dihedral type.}
There are two double cosets $P\backslash G_3/P$, represented by
$1$ and the permutation $w$ interchanging the second and third
basis vectors.  Mackey's intertwining formula gives one for the
identity-coset contribution to
$\dim\operatorname{End}_{G_3}\tau_\theta$.
For the other coset, the intersection $P\cap wPw^{-1}$ consists
of the matrices
\[
 \begin{pmatrix}a&b&c\\0&d&0\\0&0&f\end{pmatrix},
                         \qquad a d f\ne0.               \tag{G72C.10}
\]
On the subgroup $1+bE_{12}$, one of the two representations
being intertwined is the nontrivial unipotent restriction of
$C_\theta$, whereas the other is trivial.  By (G72C.8) there is
no intertwiner.  Thus the second contribution is zero and
$\dim\operatorname{End}_{G_3}\tau_\theta=1$.
Characteristic-zero semisimplicity proves irreducibility.
Finally $[G_3:P]=q^2+q+1$, which gives the degree in the statement.
\end{proof}

\begin{remark}[Scope of the character calculation]
The proof constructs an actual irreducible finite representation
and its modular character.  It does not identify a vector in any
global automorphic space.  In particular the projective lattice
$\operatorname{Ind}_{B}^{G_3}\mathcal O$ has
character $1+2\rho_{21}+\operatorname{St}_3$ and contains no
$\tau_\theta$: its projectivity follows from $p\nmid |B|$.
Thus the presence of a finite Steinberg constituent by itself
cannot replace the strictly positive multiplicity difference
required in AC-G72.
\end{remark}
